\section{ReAct：交错 Reasoning 与 Acting}

\subsection{为什么只思考或只行动都不够}

Chain-of-thought 能分解问题，但内部推理无法自动纠正过时知识或幻觉；纯工具调用能取得新信息，却缺少长期计划和跨观察整合。ReAct 让模型交替生成：

\begin{itemize}
    \item \textbf{Thought}：当前判断、子目标和下一步理由；
    \item \textbf{Action}：搜索、点击、查询或环境操作；
    \item \textbf{Observation}：工具返回的外部状态。
\end{itemize}

\videofigure{figures/react-contribution.jpg}{ReAct 用统一轨迹交错语言推理与环境动作。}{00:04:40--00:06:46}

\videofigure{figures/react-loop.jpg}{Reasoning 负责计划与整合，Action 提供新证据并校正推理。}{00:06:46--00:07:32}

\subsection{形式化为序列决策过程}

在时间步 $t$，Agent 根据历史上下文 $c_t$ 选择动作：

$$
a_t\sim\pi_\theta(a\mid c_t),\qquad
c_{t+1}=c_t\oplus a_t\oplus o_{t+1},
$$

其中 $o_{t+1}$ 是环境执行动作后返回的 observation。ReAct 的特殊之处是 $a_t$ 可包含语言 thought 和可执行 action，两者共享同一个上下文。

\videofigure{figures/react-setup.jpg}{Agent 观察环境、选择动作并把新 observation 写回上下文。}{00:07:32--00:09:06}

\begin{knowledgebox}{Thought 也是一种内部状态压缩}
模型不需要把全部历史原样保留为结构化程序状态，而是用自然语言总结当前假设、证据与待办项。这很灵活，但也会产生状态漂移：早期错误总结可能在后续上下文中被当作事实。
\end{knowledgebox}

\subsection{Reason-only、Act-only 与 ReAct}

\videofigure{figures/react-example.jpg}{在多跳问答中，ReAct 通过搜索 observation 修正内部推理。}{00:09:06--00:17:18}

\begin{itemize}
    \item \textbf{Reason-only}：能形成计划，但可能基于错误记忆持续推理；
    \item \textbf{Act-only}：能获得事实，却容易缺少子目标管理和跨步骤整合；
    \item \textbf{ReAct}：用 thought 决定查什么，用 observation 验证并更新 thought。
\end{itemize}

\subsection{知识密集任务：HotpotQA 与 FEVER}

HotpotQA 要求跨多个 Wikipedia 页面完成多跳问答；FEVER 要求检索证据判断声明真假。它们测试的不只是最终答案，还包括能否找到支撑证据。

\videofigure{figures/knowledge-tasks.jpg}{ReAct 在多跳问答与事实验证任务上结合检索和推理。}{00:17:18--00:20:46}

ReAct 的优势是轨迹可解释：研究者能看到模型查了什么、在哪一步接受了错误证据、是否重复搜索。但“可见 thought”并不保证真实反映模型内部因果过程，它更适合作为可调试工作记忆。

\subsection{决策任务：WebShop}

WebShop 要求 Agent 根据用户约束在模拟购物网站中导航、筛选并购买商品。成功不仅依赖知识，还依赖状态化操作与不可跳过的页面流程。

\videofigure{figures/webshop.jpg}{WebShop 将自然语言约束转化为多步网页行动。}{00:20:46--00:22:10}

\subsection{优势与失败模式}

\videofigure{figures/react-strengths.jpg}{ReAct 的优势是简单、可解释且无需更新模型参数。}{00:22:10--00:23:22}

常见失败包括：

\begin{itemize}
    \item 搜索查询过窄，第一条错误证据锁定后续方向；
    \item observation 太长，关键信息被噪声淹没；
    \item thought 重复描述状态但没有推进计划；
    \item 工具错误或页面变化导致轨迹失效；
    \item 缺少停止条件，Agent 在同一动作上循环。
\end{itemize}

\begin{warningbox}{工具结果不是天然可信 ground truth}
搜索结果可能过时、网页可能含恶意 prompt injection、数据库可能权限不足。Agent 必须记录来源、验证证据并限制工具权限，而不能把 observation 无条件提升为系统指令。
\end{warningbox}

\subsection{本章小结}

ReAct 把语言推理从封闭生成过程变成可与环境交互的轨迹。它主要提升当前任务中的自我修正能力，但参数并未更新；要让模型从大量执行经历中长期改进，需要训练时反馈。

